\section{Legal Landscape}
\label{sec:legal}

\subsection{Origin and Academic Status}

The Mean-Median Difference was introduced by \citet{wang2016} in a Stanford Law Review
article arguing that simple distributional statistics can identify partisan gerrymanders
with less theoretical baggage than the Efficiency Gap. Wang proposed MM as a complement
to EG, noting that its election-cycle stability made it more reliable for single-election
analysis and its independence from the swing model made it methodologically simpler to
defend against technical cross-examination.

Unlike EG, no specific MM threshold has been proposed in the academic literature as
a presumptive unconstitutionality standard. This is both a limitation (courts have
no threshold to apply) and an advantage (expert witnesses are not forced to defend
a specific threshold that opposing experts can challenge as arbitrary). The appropriate
framing in state court proceedings is comparative: not ``MM above $X\%$ is unconstitutional''
but ``enacted MM is $Y$ units above what a neutral algorithm produces for the same state,
a gap that reflects mapmaker choices rather than geography.''

\subsection{State Court Appearances}

\textbf{North Carolina} (\textit{Harper v.\ Hall}, 2022): MM was among the metrics
cited in expert testimony that the NC Supreme Court credited in invalidating the
2022 congressional map. The court treated convergent evidence across multiple metrics
as demonstrating systematic partisan intent; MM's consistency with EG and Partisan
Bias evidence strengthened the multi-metric case.

\textbf{Pennsylvania}: Expert testimony in Pennsylvania redistricting proceedings
has used MM alongside EG. PA courts have focused more on legislative intent evidence
and on the \textit{League of Women Voters} precedent's comparative analysis of
proposed maps than on specific metric thresholds.

\textbf{Wisconsin}: Wisconsin redistricting expert testimony has cited MM in state
and federal proceedings. The metric's stability advantage over EG makes it
particularly relevant in Wisconsin, where competitive statewide elections make
wave-election sensitivity a material concern.

\subsection{The Geographic Sorting Defense and Its Rebuttal}

The primary defense against MM-based gerrymandering claims is the geographic
sorting argument: positive MM reflects where Democrats choose to live, not
deliberate manipulation. This defense has significant merit as a descriptive
claim---the geographic sorting baseline is real and must be acknowledged.

The \textsc{bisect} reference distribution provides the empirical rebuttal:
\begin{enumerate}
\item Run \textsc{bisect} on the challenged state's census geography without
  partisan data, producing a map with MM $\approx 0.02$--$0.03$ (the geographic baseline).
\item Compare to the enacted map's MM $\approx 0.07$.
\item The gap ($\Delta\text{MM} \approx 0.05$) cannot be explained by geography,
  because geography is held constant between the two maps.
\item Therefore the gap reflects deliberate mapmaker choices---deliberate packing
  beyond what compactness-constrained redistricting requires.
\end{enumerate}

This rebuttal holds geography constant, which is exactly what the geographic sorting
defense requires: the defense acknowledges geographic constraints, so any MM above
the geographic baseline (measured by \textsc{bisect}) must be explained by something
other than geography.

\subsection{Recommended Expert Witness Protocol}

\begin{enumerate}
\item Report MM computed directly from observed vote shares (no swing model) for
  at least the 2016, 2018, and 2020 presidential cycles, with the standard deviation.
\item Report the \textsc{bisect} reference MM ($\approx 0.02$--$0.03$ for most states)
  as the geographic sorting baseline.
\item Report $\Delta\text{MM}$ as the mapmaker contribution, explicitly distinguishing
  it from the geographic baseline.
\item Demonstrate MM's invariance to uniform swings (Table~\ref{tab:mm-swing}) to
  preempt the wave-election sensitivity objection.
\item Note that no court has adopted a specific MM threshold; the relevant legal
  question is whether $\Delta\text{MM}$ is consistent with the applicable state
  constitutional standard for partisan gerrymandering.
\end{enumerate}
